\documentclass[11pt]{ctexart}
\usepackage[a4paper,margin=1in]{geometry}
\usepackage{longtable,booktabs}
\title{Geant4-Agent 回归测试报告（中文）}
\author{自动化评测流程}
\date{2026-03-07\_113802}
\begin{document}
\maketitle
\section{统一 Pass/Fail 总表}
\begin{longtable}{p{0.06\linewidth}p{0.08\linewidth}p{0.07\linewidth}p{0.08\linewidth}p{0.07\linewidth}p{0.07\linewidth}p{0.07\linewidth}p{0.07\linewidth}p{0.09\linewidth}p{0.06\linewidth}p{0.21\linewidth}}
\toprule
ID & 领域 & 风格 & 参数完整性 & 期望初始 & 实际初始 & 期望最终 & 实际最终 & 缺参(初->终) & 结果 & 失败原因 \\
\midrule
MTM01 & medical & fuzzy & missing & 失败 & 失败 & 通过 & 通过 & 13->0 & 通过 & - \\
MTM02 & aerospace & precise & missing & 失败 & 失败 & 通过 & 通过 & 2->0 & 通过 & - \\
MTM03 & industrial\_ndt & precise & missing & 失败 & 失败 & 通过 & 通过 & 12->0 & 通过 & - \\
MTM04 & physics & precise & missing & 失败 & 失败 & 通过 & 通过 & 8->0 & 通过 & - \\
MTM05 & nuclear\_engineering & fuzzy & missing & 失败 & 失败 & 通过 & 通过 & 13->0 & 通过 & - \\
MTM06 & security\_screening & fuzzy & missing & 失败 & 通过 & 通过 & 通过 & 0->0 & 失败 & initial\_complete\_mismatch, missing\_not\_reduced\_after\_followups \\
MTM07 & semiconductor & precise & missing & 失败 & 失败 & 通过 & 通过 & 13->0 & 通过 & - \\
MTM08 & education & fuzzy & missing & 失败 & 失败 & 通过 & 通过 & 11->0 & 通过 & - \\
MTM09 & environmental\_radiation & fuzzy & missing & 失败 & 失败 & 通过 & 通过 & 13->0 & 通过 & - \\
MTM10 & high\_energy\_physics & precise & missing & 失败 & 失败 & 通过 & 通过 & 10->0 & 通过 & - \\
\bottomrule
\end{longtable}
\section{基准对比}
\begin{itemize}
\item pass\_rate=0.900
\item initial\_alignment\_rate=0.900
\item final\_alignment\_rate=1.000
\item missing\_reduction\_rate=1.000
\item confirm\_apply\_success\_rate=1.000
\item rollback\_after\_confirm\_rate=0.000
\item delta\_pass\_rate=-0.050
\item regressions=无
\item improvements=无
\end{itemize}
\appendix
\section{附录：完整对话与内部处理链（自然语言）}
\subsection{案例 MTM01（medical）}
\textbf{对话：}
\begin{itemize}
\item 用户：???????????????????
\item 系统：你希望创建哪种几何结构呢？另外，这个结构的 module\_x 尺寸是多少？
\item 用户：?????? 300 mm x 300 mm x 300 mm??? G4\_WATER?source.type=beam?particle=gamma?energy=6 MeV?position=(0,0,-400) mm?direction=+z?physics=QBBC?output=root?
\item 系统：配置已完成。
\end{itemize}
\textbf{内部链路（自然语言）：}
\begin{itemize}
\item 第 1 轮：动作=ask\_clarification，阶段=geometry，闭环状态=False。
\item 缺失字段：geometry.structure、geometry.params.module\_x、geometry.params.module\_y、geometry.params.module\_z、materials.volume\_material\_map、source.type、source.particle、source.energy、source.position、source.direction、physics.physics\_list、output.format、output.path。
\item LLM 使用=True；回退=无；后端=llm\_semantic\_frame+runtime\_semantic。
\item 第 2 轮：动作=finalize，阶段=materials，闭环状态=True。
\item 缺失字段：无。
\item LLM 使用=True；回退=无；后端=llm\_slot\_frame+runtime\_semantic。
\end{itemize}
\subsection{案例 MTM02（aerospace）}
\textbf{对话：}
\begin{itemize}
\item 用户：I need a first-pass proton irradiation setup for a silicon device; keep it simple for now.
\item 系统：Geometry: updated geometry structure, geometry parameter module\_x, geometry parameter module\_y. Materials: updated materials.selected\_materials, volume-to-material binding, material selection source. Source: updated source type, particle type, source energy. Physics: updated physics list, physics.backup\_physics\_list, physics selection reasons. Output: still needs output format, output path.
\item 用户：Use a silicon box 50 mm x 50 mm x 5 mm, beam source proton 100 MeV from (0,0,-100) mm along +z, physics FTFP\_BERT, output json.
\item 系统：You're about to overwrite a value that's already been set.

Pending change:
- geometry parameter module\_x: 1000.0 → 50.0
- geometry parameter module\_y: 1000.0 → 50.0

Please reply **'confirm'** to apply this change, or **'keep existing'** to keep the current values.
\item 用户：confirm
\item 系统：You're about to overwrite a previously set value.

Pending change:
- geometry parameter module\_x: 1000.0 → 50.0
- geometry parameter module\_y: 1000.0 → 50.0

Please reply **'confirm'** to apply this change, or **'keep existing'** to keep the current values.
\end{itemize}
\textbf{内部链路（自然语言）：}
\begin{itemize}
\item 第 1 轮：动作=summarize\_progress，阶段=materials，闭环状态=False。
\item 缺失字段：output.format、output.path。
\item LLM 使用=True；回退=无；后端=llm\_slot\_frame+runtime\_semantic。
\item 第 2 轮：动作=confirm\_overwrite，阶段=source，闭环状态=False。
\item 缺失字段：无。
\item LLM 使用=True；回退=无；后端=llm\_slot\_frame+runtime\_semantic。
\item 第 3 轮：动作=confirm\_overwrite，阶段=physics，闭环状态=True。
\item 缺失字段：无。
\item LLM 使用=True；回退=无；后端=llm\_semantic\_frame+runtime\_semantic+confirmed\_pending\_overwrite。
\end{itemize}
\subsection{案例 MTM03（industrial\_ndt）}
\textbf{对话：}
\begin{itemize}
\item 用户：???????????????? gamma ???
\item 系统：目前我们已经确定了粒子类型为 gamma，并同步了源的选择来源和依据。

接下来还需要你补充几何结构的具体类型和尺寸参数，以及为体积指定材料。
\item 用户：??????????? 40 mm??? 80 mm??? G4\_Fe?source.type=point?particle=gamma?energy=2 MeV?position=(0,0,-200) mm?direction=+z?physics=Shielding?output=csv?
\item 系统：配置已完成。
\end{itemize}
\textbf{内部链路（自然语言）：}
\begin{itemize}
\item 第 1 轮：动作=summarize\_progress，阶段=geometry，闭环状态=False。
\item 缺失字段：geometry.structure、geometry.params.module\_x、geometry.params.module\_y、geometry.params.module\_z、materials.volume\_material\_map、source.type、source.energy、source.position、source.direction、physics.physics\_list、output.format、output.path。
\item LLM 使用=True；回退=无；后端=llm\_slot\_frame+runtime\_semantic。
\item 第 2 轮：动作=finalize，阶段=materials，闭环状态=True。
\item 缺失字段：无。
\item LLM 使用=True；回退=无；后端=llm\_slot\_frame+runtime\_semantic。
\end{itemize}
\subsection{案例 MTM04（physics）}
\textbf{对话：}
\begin{itemize}
\item 用户：Gamma point source 1 MeV at (0,0,-100) toward +z; please build the rest.
\item 系统：Great, we've set up your source: a 1 MeV gamma point source at (0,0,-100) directed along +Z.

To finish the configuration, we still need:
*   **Geometry:** The structure type and its X dimension.
*   **Materials:** Which material to assign to the volume.
*   **Physics:** The physics list to use.
*   **Output:** The format and file path for your results.
\item 用户：Set copper box 1 m x 1 m x 1 m, physics FTFP\_BERT, output hdf5.
\item 系统：Your configuration is complete.
\end{itemize}
\textbf{内部链路（自然语言）：}
\begin{itemize}
\item 第 1 轮：动作=summarize\_progress，阶段=geometry，闭环状态=False。
\item 缺失字段：geometry.structure、geometry.params.module\_x、geometry.params.module\_y、geometry.params.module\_z、materials.volume\_material\_map、physics.physics\_list、output.format、output.path。
\item LLM 使用=True；回退=无；后端=llm\_slot\_frame+runtime\_semantic。
\item 第 2 轮：动作=finalize，阶段=materials，闭环状态=True。
\item 缺失字段：无。
\item LLM 使用=True；回退=无；后端=llm\_slot\_frame+runtime\_semantic。
\end{itemize}
\subsection{案例 MTM05（nuclear\_engineering）}
\textbf{对话：}
\begin{itemize}
\item 用户：?????????????????????
\item 系统：你希望构建哪种几何结构呢？比如一个简单的盒子，还是一个圆柱体？
\item 用户：????? 200 mm x 200 mm x 200 mm??? G4\_Pb?source.type=point?particle=neutron?energy=14 MeV?position=(0,0,-50) mm?direction=+z?physics=QGSP\_BERT\_HP?output=xml?
\item 系统：配置已完成。
\end{itemize}
\textbf{内部链路（自然语言）：}
\begin{itemize}
\item 第 1 轮：动作=ask\_clarification，阶段=geometry，闭环状态=False。
\item 缺失字段：geometry.structure、geometry.params.module\_x、geometry.params.module\_y、geometry.params.module\_z、materials.volume\_material\_map、source.type、source.particle、source.energy、source.position、source.direction、physics.physics\_list、output.format、output.path。
\item LLM 使用=True；回退=无；后端=llm\_semantic\_frame+runtime\_semantic。
\item 第 2 轮：动作=finalize，阶段=materials，闭环状态=True。
\item 缺失字段：无。
\item LLM 使用=True；回退=无；后端=llm\_slot\_frame+runtime\_semantic。
\end{itemize}
\subsection{案例 MTM06（security\_screening）}
\textbf{对话：}
\begin{itemize}
\item 用户：Need a quick cargo-screening toy setup with forward photons; details can follow.
\item 系统：Your configuration is complete.
\item 用户：Use steel cylinder target radius 60 mm half-length 120 mm, beam gamma 3 MeV from (0,0,-500) mm to +z, physics Shielding, output root.
\item 系统：Your request would change two already confirmed settings:
- Geometry structure: from "single\_box" to "single\_tubs"
- Selected material: from "G4\_Cu" to "G4\_STAINLESS-STEEL"

Please reply **"confirm"** to apply these changes, or **"keep existing"** to keep the current values.
\item 用户：confirm
\item 系统：Your request changes two confirmed settings:
- Geometry structure: from "single\_box" to "single\_tubs"
- Selected material: from "G4\_Cu" to "G4\_STAINLESS-STEEL"

Please reply **"confirm"** to apply these changes, or **"keep existing"** to keep the current values.
\end{itemize}
\textbf{内部链路（自然语言）：}
\begin{itemize}
\item 第 1 轮：动作=finalize，阶段=materials，闭环状态=True。
\item 缺失字段：无。
\item LLM 使用=True；回退=无；后端=llm\_slot\_frame+runtime\_semantic。
\item 第 2 轮：动作=confirm\_overwrite，阶段=source，闭环状态=False。
\item 缺失字段：无。
\item LLM 使用=True；回退=无；后端=llm\_slot\_frame+runtime\_semantic。
\item 第 3 轮：动作=confirm\_overwrite，阶段=physics，闭环状态=True。
\item 缺失字段：无。
\item LLM 使用=True；回退=无；后端=llm\_semantic\_frame+runtime\_semantic+confirmed\_pending\_overwrite。
\end{itemize}
\subsection{案例 MTM07（semiconductor）}
\textbf{对话：}
\begin{itemize}
\item 用户：????????????????????????
\item 系统：你打算构建哪种几何结构呢？比如一个简单的盒子，还是一个圆柱体？另外，这个结构的X方向尺寸（module\_x）是多少？
\item 用户：?? 20 mm x 20 mm x 1 mm ? G4\_Si ????source.type=isotropic?particle=neutron?energy=2 MeV?position=(0,0,0) mm?physics=QGSP\_BERT\_HP?output=json?
\item 系统：配置已完成。
\end{itemize}
\textbf{内部链路（自然语言）：}
\begin{itemize}
\item 第 1 轮：动作=ask\_clarification，阶段=geometry，闭环状态=False。
\item 缺失字段：geometry.structure、geometry.params.module\_x、geometry.params.module\_y、geometry.params.module\_z、materials.volume\_material\_map、source.type、source.particle、source.energy、source.position、source.direction、physics.physics\_list、output.format、output.path。
\item LLM 使用=True；回退=无；后端=llm\_semantic\_frame+runtime\_semantic。
\item 第 2 轮：动作=finalize，阶段=materials，闭环状态=True。
\item 缺失字段：无。
\item LLM 使用=True；回退=无；后端=llm\_slot\_frame+runtime\_semantic。
\end{itemize}
\subsection{案例 MTM08（education）}
\textbf{对话：}
\begin{itemize}
\item 用户：For classroom demo, I want gamma attenuation in copper; you can ask only essential missing parts.
\item 系统：Geometry: still needs geometry structure, geometry parameter module\_x. Materials: updated materials.selected\_materials, material selection source, material selection reasons; still needs volume-to-material binding. Source: updated particle type, source selection source, source selection reasons; still needs source type, source energy. Physics: updated physics list, physics.backup\_physics\_list, physics selection reasons. Output: still needs output format, output path.
\item 用户：Okay, copper box 100 mm x 100 mm x 100 mm, point gamma 1 MeV at (0,0,-200) mm, direction +z, physics FTFP\_BERT, output csv.
\item 系统：You're about to change a confirmed setting.  
Pending change: **physics list** from *Shielding* to *FTFP\_BERT*.  

Reply **"confirm"** to apply this change, or **"keep existing"** to keep the current value.
\item 用户：confirm
\item 系统：You're about to change a confirmed setting.  
Pending change: **physics list** from *Shielding* to *FTFP\_BERT*.  

Reply **"confirm"** to apply this change, or **"keep existing"** to keep the current value.
\end{itemize}
\textbf{内部链路（自然语言）：}
\begin{itemize}
\item 第 1 轮：动作=summarize\_progress，阶段=geometry，闭环状态=False。
\item 缺失字段：geometry.structure、geometry.params.module\_x、geometry.params.module\_y、geometry.params.module\_z、materials.volume\_material\_map、source.type、source.energy、source.position、source.direction、output.format、output.path。
\item LLM 使用=True；回退=无；后端=llm\_slot\_frame+runtime\_semantic。
\item 第 2 轮：动作=confirm\_overwrite，阶段=materials，闭环状态=False。
\item 缺失字段：无。
\item LLM 使用=True；回退=无；后端=llm\_slot\_frame+runtime\_semantic。
\item 第 3 轮：动作=confirm\_overwrite，阶段=source，闭环状态=True。
\item 缺失字段：无。
\item LLM 使用=True；回退=无；后端=llm\_semantic\_frame+runtime\_semantic+confirmed\_pending\_overwrite。
\end{itemize}
\subsection{案例 MTM09（environmental\_radiation）}
\textbf{对话：}
\begin{itemize}
\item 用户：???????????????????????
\item 系统：你希望构建哪种几何结构，比如一个方盒还是一个圆柱管？另外，方盒在X方向的尺寸是多少？
\item 用户：??? 500 mm x 500 mm x 500 mm ?????? G4\_AIR?source.type=point?particle=gamma?energy=0.5 MeV?position=(0,0,-300) mm?direction=+z?physics=QBBC?output=root?
\item 系统：配置已完成。
\end{itemize}
\textbf{内部链路（自然语言）：}
\begin{itemize}
\item 第 1 轮：动作=ask\_clarification，阶段=geometry，闭环状态=False。
\item 缺失字段：geometry.structure、geometry.params.module\_x、geometry.params.module\_y、geometry.params.module\_z、materials.volume\_material\_map、source.type、source.particle、source.energy、source.position、source.direction、physics.physics\_list、output.format、output.path。
\item LLM 使用=True；回退=无；后端=llm\_semantic\_frame+runtime\_semantic。
\item 第 2 轮：动作=finalize，阶段=materials，闭环状态=True。
\item 缺失字段：无。
\item LLM 使用=True；回退=无；后端=llm\_slot\_frame+runtime\_semantic。
\end{itemize}
\subsection{案例 MTM10（high\_energy\_physics）}
\textbf{对话：}
\begin{itemize}
\item 用户：Set geometry as a sphere target; I have not decided radius or source details yet.
\item 系统：I've updated your geometry to a sphere target. To continue, I still need:

*   The sphere's maximum radius (child\_rmax).
*   To assign a material to your volume.
*   Your source setup: type, particle, and energy.
*   Your physics list selection.
*   Your output preferences: format and file path.
\item 用户：Sphere target radius 80 mm, lead material, electron point source 5 MeV at (0,0,-200) mm toward +z, physics FTFP\_BERT, output hdf5.
\item 系统：Your configuration is complete.
\end{itemize}
\textbf{内部链路（自然语言）：}
\begin{itemize}
\item 第 1 轮：动作=summarize\_progress，阶段=geometry，闭环状态=False。
\item 缺失字段：geometry.params.child\_rmax、materials.volume\_material\_map、source.type、source.particle、source.energy、source.position、source.direction、physics.physics\_list、output.format、output.path。
\item LLM 使用=True；回退=无；后端=llm\_slot\_frame+runtime\_semantic。
\item 第 2 轮：动作=finalize，阶段=materials，闭环状态=True。
\item 缺失字段：无。
\item LLM 使用=True；回退=无；后端=llm\_slot\_frame+runtime\_semantic。
\end{itemize}
\end{document}